\scanpage{141}

% L1, introduction
Nous pouvons maintenant énoncer

% L2–4, ruled theorem
\uline{Thm} 2 Supposons $\mathrm{rg}\,E=N$ pair $=2n$ \inserted{et supposons 2 inv. sur $S$.} Alors le morphisme précédent est lisse projectif à fibres \inserted{géom} $0$-connexes.

% L5–6, proof
Lisse projectif est évident, il faut voir que les fibres \inserted{géom} sont $0$-connexes.

% L7–10, cancelled proof
\deleted{Elles sont non vides : il suffit \inserted{(sur $k$ alg. clos)} de trouver \deleted{$\omega_{\unsure{F}}$} $F'$ avec $\omega_{F'}=-\omega_F$. Il suffit pour ceci de \deleted{remplacer} prendre}

% L10–11, cancelled display
\[
\begin{aligned}
\deleted{F'}&\deleted{=\deleted{\sum\unsure{\mathcal O}}\mathcal O_S e_1^*+\sum_{2\leq i\leq n}\underline O_S e_i}\\
\deleted{G'}&\deleted{=\mathcal O_S e_1+\sum_{2\leq i\leq n}\underline O_S e_i^*}
\end{aligned}
\]

% L12–17, lower text
Le fait que la flèche est surjective provient du fait que $\underline{\mathrm{Aut}}(E,q)=O(q)$ opère transitivement sur $\underline{\mathrm{Arf}}(E,q)$ (une \unsure{réflexion|reflexion} orthogonale p.ex. \deleted{\unsure{à déduire que}} \inserted{agit} comme \deleted{\unsure{symétrie fondamentale}} \inserted{l'opération canonique} du rev. d'Arf), il reste à prouver que les

% L18–20, lower continuation
\inserted{schéma} \inserted{$\underline\Delta_{\{n\}}^\omega(E,q)$} $F$ tels que $\omega_F=\omega$ (fixé) \deleted{forment} \unsure{forment|forme} un schéma connexe (sur $k$ alg. clos). Or c'est l'image de

% L21, bottom continuation
$\underline\Delta_{\deleted{\unsure{\{H\}}}\unsure{H\setminus\{1\}|\{1\}|\{1,\ldots,n\}}}^{\omega}$, où $H=[1,n]$, donc on est

% bottom left, pencil
% archive sheet number: 140
